\documentclass[../proyecto.tex]{subfiles}
\graphicspath{{\subfix{../images/}}}
\begin{document}

\chapter{Apéndice}

% Se puede incorporar imágenes de obras, links a obras sonoras o audiovisuales, entrevistas, material de archivo, etc.

\section{Relo}

Relo es un módulo electrónico que está en desarrollo en el marco de esta tesis. Permite controlar la velocidad de una señal de reloj, usada comúnmente para controlar la velocidad de secuenciadores en un contexto de sintetizadores.

\begin{figure}[htbp]
    \centering
    \includegraphics[width=0.5\textwidth]{./imagenes/relov0.png}
    \caption{Una PCB diseñada del módulo Relo}
    \label{fig:relo}
\end{figure}

\section{Parla}

Parla es un módulo electrónico en desarrollo durante esta tesis. Ha contado con el diseño gráfico de Mateo Arce, y con el apoyo en diseño de placas y estándares eléctricos de Matías Serrano.

\begin{figure}[htbp]
    \centering
    \includegraphics[width=0.5\textwidth]{imagenes/parlav0revB}
    \caption{Un esquemático de una de las iteraciones del módulo Parla}
    \label{fig:parla}
\end{figure}


\section{Maqui}

Maqui es un secuenciador que se puede usar para controlar otros módulos de sonido.

Este trabajo cuenta con la colaboración de la primera persona practicante que ha tenido la empresa Piruetas Spa, el diseñador Andrés Martin.

\section{Allender}

Allender es un sintetizador de tipo sampler, que está cargado con samples extraídos programáticamente de discursos del presidente Salvador Allende.

\section{Silenkier}

Silenkier es un sintetizador de tipo sampler, que está cargado con distintos samples de audio extraídos programáticamente de los discursos del actual presidente de Chile, que fueron percibidos como silencios por algoritmos de detección de habla.

\end{document}